\documentclass[11pt]{article}

% Answers and hints for the practice problems of Lecture 1
\usepackage[margin=1in]{geometry}
\usepackage[T1]{fontenc}
\usepackage{lmodern}
\usepackage{amsmath,amssymb,amsthm}
\usepackage{enumitem}
\usepackage{xr-hyper}
\usepackage{hyperref}
% Problem numbers are read from practice1.aux, so compile practice1.tex first.
\externaldocument{practice1}[practice1.pdf]
\usepackage{xcolor}

\title{Lecture 1: Practice Problems --- Answers and Hints}
\author{Linear Algebra}
\date{\today}

\DeclareMathOperator{\im}{im}
\DeclareMathOperator{\rank}{rank}
\DeclareMathOperator{\Null}{null}
\DeclareMathOperator{\col}{col}
\DeclareMathOperator{\row}{row}
\DeclareMathOperator{\Span}{span}
\DeclareMathOperator{\RREF}{RREF}
\DeclareMathOperator{\tr}{tr}
\DeclareMathOperator{\diag}{diag}

\begin{document}

\maketitle

\section*{Number Sets and Fields}

\begin{description}[leftmargin=0pt, itemsep=0.3em]
    \item[\ref{pr:numbers}.] $\frac{7}{3}, 0.25 \in \mathbb{Q}, \mathbb{R}, \mathbb{C}$; $\sqrt{49} = 7$, $\pi - \pi = 0$ and $(1 + i)(1 - i) = 2$ lie in all four sets; $\sqrt{12} = 2\sqrt{3} \in \mathbb{R}, \mathbb{C}$; $2 + 3i \in \mathbb{C}$ only.
    \item[\ref{pr:field-fail}.] (a) Additive inverses: $-1 \notin \mathbb{N}$. (b) There is no additive identity, since $0 \notin \mathbb{R}_{>0}$. (c) Not closed under addition: $1 + 1 = 2$. (d) There is no multiplicative identity, since $1 \notin 2\mathbb{Z}$.
    \item[\ref{pr:Qsqrt2}.] (a) $(a + b\sqrt{2})^{-1} = \frac{a}{a^2 - 2b^2} - \frac{b}{a^2 - 2b^2}\sqrt{2}$; if $a^2 = 2b^2$ with $(a, b) \neq (0, 0)$, then $\sqrt{2} = |a/b|$ would be rational. (b) No: $2$ has no inverse, since $\frac{1}{2} = a + b\sqrt{2}$ forces $b = 0$ and $a = \frac{1}{2} \notin \mathbb{Z}$.
\end{description}

\section*{Vector Spaces}

\begin{description}[leftmargin=0pt, itemsep=0.3em]
    \item[\ref{pr:weird-ops}.] (a) Only $1 \odot v = v$ fails. (b) Only $(a + b) \odot v = a \odot v + b \odot v$ fails: the second coordinates are $y$ and $2y$.
    \item[\ref{pr:positive-reals}.] The zero vector is $1$, and the additive inverse of $u$ is $\frac{1}{u}$. For example, $(a + b) \odot u = u^{a + b} = u^a u^b = (a \odot u) \oplus (b \odot u)$.
    \item[\ref{pr:functions}.] (a) The zero vector is the zero function $x \mapsto 0$. (b) No: the zero function is not in the set (and the set is not closed under addition).
    \item[\ref{pr:Cn-membership}.] (a) $n = \infty$: all derivatives of $e^{-x^2}$ are polynomials times $e^{-x^2}$. (b) $n = 2$: $|x|^3 = x^2|x|$, with $(|x|^3)' = 3x|x|$ and $(|x|^3)'' = 6|x|$, which is not differentiable at $0$ (compare the exercise on $x^n|x|$ in the lecture). (c) $n = 0$: $\sqrt{|x|}$ is continuous, but not differentiable at $0$, since $\frac{\sqrt{|h|}}{h}$ is unbounded as $h \to 0$. (d) $n = \infty$: $f(x) = 1 - \frac{x^2}{3!} + \frac{x^4}{5!} - \cdots$ is a power series that converges for all $x$, and such power series can be differentiated term by term as often as we like. (e) $g$ is not continuous at $0$.
    \item[\ref{pr:fn-sets}.] (a) Yes: the zero function has period $1$, and if $f$ and $g$ have period $1$, then so do $f + g$ and $cf$. (b) No: $f(x) = x$ is increasing, but $(-1) \cdot f$ is not. (c) Yes, by the rules for limits. (d) No: the constant function $1$ lies in the set, but $2 \cdot 1$ does not. (e) Yes.
\end{description}

\section*{Bases and Dimension}

\begin{description}[leftmargin=0pt, itemsep=0.3em]
    \item[\ref{pr:indep}.] Dependent: $(1, 2, 1) - (2, 1, 0) + (1, -1, -1) = 0$. The span has basis $\{(1, 2, 1), (2, 1, 0)\}$ and dimension $2$.
    \item[\ref{pr:basisU}.] $\dim U = 2$, with basis $(1, 0, -2, 1)$, $(0, 1, -1, 0)$.
    \item[\ref{pr:taylor-basis}.] Substitute $x = t + 1$: $[x^2]_{\mathcal{B}} = (1, 2, 1)$ and $[3 + 2x + x^2]_{\mathcal{B}} = (6, 4, 1)$.
\end{description}

\section*{Linear Transformations}

\begin{description}[leftmargin=0pt, itemsep=0.3em]
    \item[\ref{pr:linear}.] (a) Linear. (b) Not linear: $T(0, 0) \neq (0, 0)$. (c) Not linear: $T(2v) = 4T(v)$ in the first coordinate. (d) Linear. (e) Not linear: $T(-1) = 1 \neq -T(1)$. (f) Linear.
    \item[\ref{pr:kerim}.] $\ker(T) = \Span\{(1, -1, 1)\}$ and $\im(T) = \{(a, b, c) : c = a - b\} = \Span\{(1, 0, 1), (1, 1, 0)\}$. $T$ is neither injective nor surjective.
    \item[\ref{pr:gallery}.] (a) $T(\mathbf{e}_1) = \mathbf{e}_2$, $T(\mathbf{e}_2) = \mathbf{e}_1$; bijective; the F is mirrored. (b) $\ker(T)$ is the $y$-axis and $\im(T)$ is the $x$-axis; the square collapses onto the segment from $(0, 0)$ to $(1, 0)$. (c) Bijective; the square becomes the parallelogram with vertices $(0, 0), (1, 0), (3, 1), (2, 1)$. (d) Rotation by $180^\circ$; bijective; the F is not mirrored. Only (a) mirrors the F.
    \item[\ref{pr:derivative}.] $\ker(D)$ is the set of constant polynomials, and $\im(D) = \mathbb{R}[x]_{\leq 2}$. $D$ is neither injective nor surjective. $D \circ D \circ D \circ D = 0$.
    \item[\ref{pr:composition}.] (a) If $f(a) = f(b)$, then $g(f(a)) = g(f(b))$, so $a = b$. (b) $f: \{1\} \to \{1, 2\}$, $f(1) = 1$, and $g: \{1, 2\} \to \{1\}$. (d) $A$ is injective and $B$ is surjective.
    \item[\ref{pr:distributive}.] (a) For example, $f(x) = x + 1$ and $g(x) = 2x$. (c) With $f(x) = x^2$ and $g = h = \mathrm{id}$: $f(2x) = 4x^2 \neq 2x^2$. (d) $c(g(x) + h(x)) = cg(x) + ch(x)$.
    \item[\ref{pr:iso}.] (a) $T(a_0 + a_1 x + a_2 x^2 + a_3 x^3) = (a_0, a_1, a_2, a_3)$. (b) $\log: V \to \mathbb{R}$, since $\log(uv) = \log u + \log v$ and $\log(u^c) = c\log u$. (c) Yes; no (dimensions $2$ and $3$); yes.
    \item[\ref{pr:calculus-maps}.] (a), (b) and (d) are linear: (a) by the sum and constant factor rules for derivatives, (b) by the linearity of the integral, and (d) since $(f + g)(x + 1) = f(x + 1) + g(x + 1)$ and $(cf)(x + 1) = c f(x + 1)$. (c) is not linear: replacing $f$ by $2f$ multiplies the value by $4$. (e) is not linear, since it maps the zero function to the constant function $1$, but a linear transformation maps $0$ to $0$. (f) is not linear: $2f$ is mapped to $4f(0)f(1)$.
    \item[\ref{pr:diff-int}.] (a) $D(J\cos) = \cos$ and $J(D\cos)(x) = \int_0^x (-\sin t)\, dt = \cos x - 1$; $D(Je^{2x}) = e^{2x}$ and $J(De^{2x})(x) = \int_0^x 2e^{2t}\, dt = e^{2x} - 1$. (b) By the fundamental theorem of calculus, $J(Df)(x) = \int_0^x f'(t)\, dt = f(x) - f(0)$. So $J(Df) = f$ exactly for the functions with $f(0) = 0$. (c) $D^2$ is a composition of linear transformations, hence linear; $\cos'' = -\cos$ and $\sin'' = -\sin$.
\end{description}

\section*{Subspaces}

\begin{description}[leftmargin=0pt, itemsep=0.3em]
    \item[\ref{pr:subspaces}.] (a) yes; (b) no; (c) no; (d) no; (e) yes; (f) yes.
    \item[\ref{pr:intersection}.] (b) The $x$-axis and the $y$-axis in $\mathbb{R}^2$: $(1, 0) + (0, 1) = (1, 1)$ lies in neither. (c) If $u \in U \setminus W$ and $w \in W \setminus U$, then $u + w$ lies in neither $U$ nor $W$.
    \item[\ref{pr:poly-subspace}.] It is the kernel of $p \mapsto p(1)$. A basis is $x - 1, \ x^2 - 1, \ x^3 - 1$, so the dimension is $3$.
    \item[\ref{pr:fn-subspaces}.] (a) $U$ is the kernel of the linear transformation $C^2(\mathbb{R}) \to C(\mathbb{R})$, $f \mapsto f'' + f$, so it is a subspace. Since $\cos'' = -\cos$ and $\sin'' = -\sin$, we get $(a\cos x + b\sin x)'' = -(a\cos x + b\sin x)$. (b) The zero function is both even and odd, and sums and multiples of even (odd) functions are again even (odd). In the hint, the first term is even and the second one is odd. (c) Yes: it is the kernel of the linear transformation $f \mapsto f(0) - f'(0)$. No: it does not contain the zero polynomial (and $x^2 + (x - x^2) = x$ has degree $1$).
\end{description}

\section*{Matrices}

\begin{description}[leftmargin=0pt, itemsep=0.3em]
    \item[\ref{pr:matvec}.] (a) $A\mathbf{x} = \mathbf{0}$. (b) $AA^\top = \begin{pmatrix} 5 & 1 \\ 1 & 11 \end{pmatrix}$ and $A^\top A = \begin{pmatrix} 2 & -3 & 1 \\ -3 & 9 & 3 \\ 1 & 3 & 5 \end{pmatrix}$, both symmetric.
    \item[\ref{pr:sizes}.] $AB$: $2 \times 3$; $BA$: undefined; $AC$: $2 \times 2$; $CA$: $3 \times 3$; $A + C^\top$: $2 \times 3$; $A^\top + C$: $3 \times 2$; $ABC$: $2 \times 2$.
    \item[\ref{pr:reflect-rotate}.] $FR = \begin{pmatrix} 0 & -1 \\ -1 & 0 \end{pmatrix}$ is the reflection in the line $y = -x$, and $RF = \begin{pmatrix} 0 & 1 \\ 1 & 0 \end{pmatrix}$ is the reflection in the line $y = x$. So $FR \neq RF$.
    \item[\ref{pr:rotation-addition}.] (b) $\cos(\alpha + \beta) = \cos\alpha\cos\beta - \sin\alpha\sin\beta$ and $\sin(\alpha + \beta) = \sin\alpha\cos\beta + \cos\alpha\sin\beta$.
    \item[\ref{pr:complex-matrix}.] (c) $M(1 + i)^{10} = 2^5 R_{10\pi/4} = 32 R_{\pi/2} = \begin{pmatrix} 0 & -32 \\ 32 & 0 \end{pmatrix} = M(32i)$, so $(1 + i)^{10} = 32i$.
    \item[\ref{pr:ref}.] The first is in RREF. The second is in REF, but not RREF. The third and fourth are not in REF (the pivots move left; a zero row lies above a nonzero row). The fifth is not in REF with our convention, since its first pivot is $2$.
    \item[\ref{pr:rref-compute}.] $\RREF(P) = \begin{pmatrix} 1 & 2 & 0 & 1 \\ 0 & 0 & 1 & -1 \\ 0 & 0 & 0 & 0 \end{pmatrix}$, and $\rank(P) = 2$.
    \item[\ref{pr:system}.] The RREF has rows $(1, 0, 5, 8)$ and $(0, 1, -1, -2)$, so the solutions are
    \[(w, x, y, z) = y\,(-5, 1, 1, 0) + z\,(-8, 2, 0, 1), \qquad y, z \in \mathbb{R}.\]
    The solution $(0, 2, -8, 5)$ is the case $y = -8$, $z = 5$; it is not a multiple of $(-5, 1, 1, 0)$. The dimension is $2$.
    \item[\ref{pr:inverse2}.] $A^{-1} = \begin{pmatrix} 2 & -1 \\ -5 & 3 \end{pmatrix}$.
    \item[\ref{pr:inverse3}.] $A^{-1} = \begin{pmatrix} 0 & -1 & 1 \\ 1 & 0 & -1 \\ -1 & 1 & 1 \end{pmatrix}$ and $\mathbf{x} = (1, -2, 4)^\top$.
    \item[\ref{pr:arithmetic}.] (a) $\mathbf{x} = (1, -2, 1)^\top$. (b) row $1$ $-$ $2\cdot$row $2$ $+$ row $3$ $= 0$, so row operations produce a zero row.
    \item[\ref{pr:elementary}.] Swap the rows, multiply row $2$ by $\frac{1}{2}$, subtract $3$ times row $2$ from row $1$: $E_1 = \begin{pmatrix} 0 & 1 \\ 1 & 0 \end{pmatrix}$, $E_2 = \begin{pmatrix} 1 & 0 \\ 0 & \frac{1}{2} \end{pmatrix}$, $E_3 = \begin{pmatrix} 1 & -3 \\ 0 & 1 \end{pmatrix}$. Then $A^{-1} = E_3 E_2 E_1 = \begin{pmatrix} -\frac{3}{2} & 1 \\ \frac{1}{2} & 0 \end{pmatrix}$ and $A = E_1^{-1} E_2^{-1} E_3^{-1}$.
    \item[\ref{pr:idempotent}.] (a) $A = A^{-1}A^2 = A^{-1}A = I$. (b) $\ker(T_P) = \Span\{(1, -1)\}$ and $\im(T_P)$ is the $x$-axis: $T_P$ projects onto the $x$-axis along lines parallel to $(1, -1)$. (d) If $\mathbf{v} = A\mathbf{w}$ and $A\mathbf{v} = \mathbf{0}$, then $\mathbf{v} = A\mathbf{w} = A^2 \mathbf{w} = A\mathbf{v} = \mathbf{0}$.
    \item[\ref{pr:nilpotent}.] $N^2 = \begin{pmatrix} 0 & 0 & 3 \\ 0 & 0 & 0 \\ 0 & 0 & 0 \end{pmatrix}$, $N^3 = 0$, and $(I - N)^{-1} = I + N + N^2 = \begin{pmatrix} 1 & 1 & 5 \\ 0 & 1 & 3 \\ 0 & 0 & 1 \end{pmatrix}$.
    \item[\ref{pr:interpolation}.] (a) $\begin{pmatrix} 1 & -1 & 1 \\ 1 & 0 & 0 \\ 1 & 1 & 1 \end{pmatrix}$, with determinant $2$. (c) $p = 1 - x^2$ and $q = 2 + 2x + x^2$.
    \item[\ref{pr:dims-special}.] (a) $3$, basis $e_{11}, e_{22}, e_{12} + e_{21}$; (b) $1$, basis $e_{12} - e_{21}$; (c) $6$; (d) $10$; (e) $6$; (f) $5$.
    \item[\ref{pr:groups}.] (a) Yes. (b) No: $\begin{pmatrix} 1 & 0 \\ 0 & 2 \end{pmatrix}\begin{pmatrix} 0 & 1 \\ 1 & 0 \end{pmatrix} = \begin{pmatrix} 0 & 1 \\ 2 & 0 \end{pmatrix}$ is not symmetric. (c) No: $\begin{pmatrix} 1 & 0 \\ 0 & -1 \end{pmatrix}$ and $\begin{pmatrix} 0 & 1 \\ 1 & 0 \end{pmatrix}$ square to $I_2$, but their product squares to $-I_2$. (d) Yes. (e) Yes; there are $n!$ of them, and each one is orthogonal.
    \item[\ref{pr:inverse-triangular}.] $\begin{pmatrix} 1 & -2 & 5 \\ 0 & 1 & -4 \\ 0 & 0 & 1 \end{pmatrix}$. For an upper triangular $A$ with nonzero diagonal entries, the Gauss--Jordan method only needs to scale rows and to subtract multiples of lower rows from higher rows, and these operations keep the right block upper triangular.
    \item[\ref{pr:orthogonal2}.] $Q = R_\theta$ (rotations, determinant $1$) or $Q = R_\theta \begin{pmatrix} 1 & 0 \\ 0 & -1 \end{pmatrix} = \begin{pmatrix} \cos\theta & \sin\theta \\ \sin\theta & -\cos\theta \end{pmatrix}$ (reflections, determinant $-1$).
    \item[\ref{pr:tridiagonal}.] (a) $D_1 = 2$, $D_2 = 3$, $D_3 = 4$. (c) $D_n = n + 1$.
    \item[\ref{pr:det4}.] $\det(A) = -1$, so $A$ is invertible.
    \item[\ref{pr:det-rules}.] $-6$, $18$, $54$, $-\frac{3}{2}$, $2$; and $\det(C) = -2$.
    \item[\ref{pr:vandermonde}.] (a) Subtract row $1$ from rows $2$ and $3$, and take out the factors $b - a$ and $c - a$. (b) The determinant is nonzero, so the map $p \mapsto (p(a), p(b), p(c))$ is an isomorphism.
    \item[\ref{pr:lambda}.] (a) $\lambda = 2$ and $\lambda = 5$; $\det(A - 3I_3) = -16 \neq 0$. (b) $\mathbf{x} = (-8, 1, 1)$ for $\lambda = 2$, and $\mathbf{x} = (-1, -1, 1)$ for $\lambda = 5$. (c) Every row sums to $11$, so $A(1, 1, 1)^\top = 11 (1, 1, 1)^\top$.
    \item[\ref{pr:upper5}.] (a) $\det(A) = 1 \cdot 2 \cdot 3 \cdot 4 \cdot 5 = 120$; yes; yes. (b) $(0, 0, 0, 0, 25)$. (c) $A - I_5$ and $A - 2I_5$ are triangular with a $0$ on the diagonal, so their determinants are $0$. Take $\mathbf{x} = \mathbf{e}_1$ and $\mathbf{y} = (2, 1, 0, 0, 0)^\top$.
    \item[\ref{pr:area}.] (a) $\frac{1}{2}\left|\det \begin{pmatrix} 3 & 1 \\ 1 & 4 \end{pmatrix}\right| = \frac{11}{2}$. (b) $\det = 7$: the areas are $7$ and $7\pi$. (d) Only the reflection, with determinant $-1$. The rotation by $180^\circ$ has determinant $+1$.
    \item[\ref{pr:fibonacci}.] $M^{10} = \begin{pmatrix} 89 & 55 \\ 55 & 34 \end{pmatrix}$, and $89 \cdot 34 - 55^2 = 1$. Cassini follows from $\det(M^n) = \det(M)^n = (-1)^n$.
    \item[\ref{pr:trace}.] (a) For example, $A = \begin{pmatrix} 1 & 1 \\ 0 & 1 \end{pmatrix}$ and $B = \begin{pmatrix} 0 & 1 \\ 1 & 0 \end{pmatrix}$. $\tr(AB - BA) = 0 \neq 2 = \tr(I_2)$. (b) $\tr(A^\top A)$ is the sum of the squares of all entries. (c) $\tr(A\mathbf{e}_{ji}) = A_{ij}$.
    \item[\ref{pr:01-matrices}.] (a) $6$. (For $3 \times 3$ matrices with entries in $\{0, 1\}$, a computer finds that $174$ of the $512$ are invertible.) (b) If $A$ and $A^{-1}$ have integer entries, then $\det(A)\det(A^{-1}) = 1$ is a product of two integers. Conversely, use the cofactor formula $A^{-1} = \frac{1}{\det(A)} C^\top$.
    \item[\ref{pr:lu-compute}.] (a) The multipliers are $3$ and $0$ in the first column and $2$ in the second: $L = \begin{pmatrix} 1 & 0 & 0 \\ 3 & 1 & 0 \\ 0 & 2 & 1 \end{pmatrix}$, $U = \begin{pmatrix} 1 & 2 & 1 \\ 0 & 2 & -2 \\ 0 & 0 & 5 \end{pmatrix}$. (b) $\det(A) = 1 \cdot 2 \cdot 5 = 10$. Forward substitution gives $\mathbf{y} = (2, 4, -5)$, and back substitution gives $\mathbf{x} = (1, 1, -1)$.
    \item[\ref{pr:lu-pivot}.] $LU = \begin{pmatrix} u_{11} & u_{12} \\ l u_{11} & lu_{12} + u_{22} \end{pmatrix}$ would need $u_{11} = 0$ and $lu_{11} = 3$. With $P = \begin{pmatrix} 0 & 1 \\ 1 & 0 \end{pmatrix}$, $PA = \begin{pmatrix} 3 & 1 \\ 0 & 2 \end{pmatrix}$ is already upper triangular: $L = I_2$ and $U = PA$.
    \item[\ref{pr:schur-practice}.] (a) $S = \begin{pmatrix} 2 & 1 \\ 1 & 2 \end{pmatrix} - \frac{1}{2}\begin{pmatrix} 1 & 1 \\ 1 & 1 \end{pmatrix} = \frac{1}{2}\begin{pmatrix} 3 & 1 \\ 1 & 3 \end{pmatrix}$, $\det(S) = 2$, $\det(M) = 4$. (b) $S^{-1} = \frac{1}{4}\begin{pmatrix} 3 & -1 \\ -1 & 3 \end{pmatrix}$, and $M^{-1} = \frac{1}{4}\begin{pmatrix} 3 & -1 & -1 \\ -1 & 3 & -1 \\ -1 & -1 & 3 \end{pmatrix}$.
    \item[\ref{pr:sherman-morrison-practice}.] $A^{-1}\mathbf{u} = (1, \frac{1}{2}, \frac{1}{4})^\top$ and $1 + \mathbf{v}^\top A^{-1}\mathbf{u} = \frac{11}{4}$, so $\det(A + \mathbf{u}\mathbf{v}^\top) = 8 \cdot \frac{11}{4} = 22$ and $(A + \mathbf{u}\mathbf{v}^\top)^{-1} = \frac{1}{22}\begin{pmatrix} 14 & -4 & -2 \\ -4 & 9 & -1 \\ -2 & -1 & 5 \end{pmatrix}$.
    \item[\ref{pr:det-lemma-practice}.] (a) $\det(aI_n + b\mathbf{1}\mathbf{1}^\top) = a^n\det(I_n + \frac{b}{a}\mathbf{1}\mathbf{1}^\top) = a^n(1 + \frac{nb}{a})$. (b) The matrix is $(\mathbf{x} - 1)I_n + J_n$, with determinant $(\mathbf{x} - 1)^{n-1}(\mathbf{x} - 1 + n)$ (for $\mathbf{x} = 1$, it is $J_n$, which is singular for $n \geq 2$). It is singular exactly for $\mathbf{x} = 1$ and $\mathbf{x} = 1 - n$ (if $n \geq 2$). (c) $\frac{1}{a}I_n - \frac{b}{a(a + nb)}J_n$.
    \item[\ref{pr:push-through}.] (a) $X^\top(XX^\top + \lambda I_N) = X^\top XX^\top + \lambda X^\top = (X^\top X + \lambda I_d)X^\top$; multiply by $(X^\top X + \lambda I_d)^{-1}$ from the left and by $(XX^\top + \lambda I_N)^{-1}$ from the right (both matrices are invertible, as in the lecture). (b) Here $XX^\top + \lambda I_1 = \mathbf{x}^\top \mathbf{x} + \lambda$ is a number. Directly: $(\mathbf{x}\mathbf{x}^\top + \lambda I_d)\mathbf{x} = \mathbf{x}(\mathbf{x}^\top \mathbf{x}) + \lambda \mathbf{x} = (\mathbf{x}^\top \mathbf{x} + \lambda)\mathbf{x}$.
\end{description}

\section*{Null Space, Column Space and Rank}

\begin{description}[leftmargin=0pt, itemsep=0.3em]
    \item[\ref{pr:bases-B}.] $\RREF(B)$ has rows $(1, 2, 0, 3)$, $(0, 0, 1, 2)$, $(0, 0, 0, 0)$, and $\rank(B) = 2$. Bases: $\Null(B)$: $(-2, 1, 0, 0)$, $(-3, 0, -2, 1)$; $\col(B)$: $(1, 2, 1)$, $(0, 1, 1)$; $\row(B)$: $(1, 2, 0, 3)$, $(0, 0, 1, 2)$. Indeed, $2 + 2 = 4$.
    \item[\ref{pr:onto}.] (a) $5 - 3 = 2$. (b) Not surjective, since the rank is at most $3$; it can be injective. (c) No: $\dim \Null(A) \geq 6 - 4 = 2$. (d) $3$.
    \item[\ref{pr:intersect}.] (a) The $7 \times 8$ matrix has a nonzero null vector $(a_1, \ldots, a_4, b_1, \ldots, b_4)$; then $\sum a_i \mathbf{v}_i = -\sum b_j \mathbf{w}_j$ is a nonzero vector in $V \cap W$. (c) $\Span\{\mathbf{e}_1, \mathbf{e}_2, \mathbf{e}_3\}$ and $\Span\{\mathbf{e}_4, \mathbf{e}_5, \mathbf{e}_6\}$.
    \item[\ref{pr:rankAB}.] (b) $\rank(AB) \leq \rank(A) \leq 2 < 3$.
    \item[\ref{pr:ATA}.] (a) $\mathbf{x}^\top A^\top A \mathbf{x} = \sum_i (A\mathbf{x})_i^2$. (b) Use the rank--nullity theorem. (c) $\Null(A^\top A) = \Null(A) = \{\mathbf{0}\}$.
\end{description}

\section*{Dual Vector Spaces}

\begin{description}[leftmargin=0pt, itemsep=0.3em]
    \item[\ref{pr:functionals}.] (a) Yes, $a = (2, 0, -1)$. (b) No. (c) Yes. (d) No: $\det(2A) = 4\det(A)$. (e) Yes. (f) No.
    \item[\ref{pr:dual-basis-r2}.] (a) $P = \begin{pmatrix} 2 & 1 \\ 1 & 1 \end{pmatrix}$ has the inverse $\begin{pmatrix} 1 & -1 \\ -1 & 2 \end{pmatrix}$, so $\mathbf{b}^1(\mathbf{x}) = x_1 - x_2$ and $\mathbf{b}^2(\mathbf{x}) = -x_1 + 2x_2$. (b) $\varphi = \varphi(b_1) \mathbf{b}^1 + \varphi(b_2) \mathbf{b}^2 = 7\mathbf{b}^1 + 4\mathbf{b}^2$.
    \item[\ref{pr:lagrange-simpson}.] (a) $L_0 = \frac{(x - 1)(x - 2)}{2}$, $L_1 = -x(x - 2)$, $L_2 = \frac{x(x - 1)}{2}$. (b) $p = L_0 + 2L_1 + 5L_2$, since $p(0) = 1$, $p(1) = 2$, $p(2) = 5$. (c) $I(L_0) = \frac{1}{3}$, $I(L_1) = \frac{4}{3}$, $I(L_2) = \frac{1}{3}$, so $\int_0^2 p(x)\, dx = \frac{1}{3}\big(p(0) + 4p(1) + p(2)\big)$.
    \item[\ref{pr:dual-map}.] (a) $(2, -1)A = (2, -1)\begin{pmatrix} 1 & 1 & 0 \\ 0 & 1 & -1 \end{pmatrix} = (2, 1, 1)$. (b) $\psi(T(\mathbf{v})) = \psi(3, -1) = 7$ and $(2, 1, 1)\mathbf{v} = 7$. (c) $A^\top$.
    \item[\ref{pr:annihilator}.] (a) The rows $\mathbf{a}$ with $a_1 + 2a_2 = 0$ and $a_2 + a_3 = 0$, spanned by $(2, -1, 1)$; $2 + 1 = 3$. (b) $(1, 0, -1, 0)$ and $(0, 1, 0, -1)$.
    \item[\ref{pr:certificate}.] (a) $\mathbf{y} = (2, -1)$: $\mathbf{y}A = \mathbf{0}$ and $\mathbf{y}\mathbf{b} = -1$. (b) The rows with $\mathbf{y}A = \mathbf{0}$ are those with $y_1 + 2y_2 + 3y_3 = 0$, spanned by $(-2, 1, 0)$ and $(-3, 0, 1)$; so the system is solvable if and only if $b_2 = 2b_1$ and $b_3 = 3b_1$, i.e., $\mathbf{b}$ is a multiple of $(1, 2, 3)^\top$.
    \item[\ref{pr:double-dual}.] $\iota(\mathbf{e}_1)(\varphi) = \varphi(\mathbf{e}_1) = a_1$: the functional ``first entry of the row''. If $\iota(\mathbf{v}) = 0$, then $\varphi(\mathbf{v}) = 0$ for all rows, in particular for $\mathbf{e}_1^\top$ and $\mathbf{e}_2^\top$, so $v_1 = v_2 = 0$.
\end{description}

\section*{Quotient Spaces and the Isomorphism Theorems}

\begin{description}[leftmargin=0pt, itemsep=0.3em]
    \item[\ref{pr:cosets}.] (a) Yes, since $(2, 1)^\top \in U$; yes, since $(4, 2)^\top \in U$; no, since $(1, 2)^\top \notin U$. (b) $\mathbf{v} - \mathbf{w} \in U$ if and only if $v_1 - 2v_2 = w_1 - 2w_2$; the representative on the $x_1$-axis is $(c, 0)^\top$. (c) $x + U \mapsto x_1 - 2x_2$.
    \item[\ref{pr:quotient-dim}.] (a) $3$, so $\mathbb{R}^5/U \cong \mathbb{R}^3$. (b) $9 - 6 = 3$; it is isomorphic to $\mathrm{Skew}_3$ (Problem~\ref{pr:first-iso}). (c) The map $p \mapsto (p(1), p(2))$ is surjective with kernel $U$, so $\mathbb{R}[x]_{\leq 4}/U \cong \mathbb{R}^2$ and $\dim U = 3$.
    \item[\ref{pr:first-iso}.] (a) $T$ is surjective with kernel $P$, so $\mathbb{R}^3/P \cong \mathbb{R}$. (b) $p \mapsto (p(0), p(1))$ is surjective onto $\mathbb{R}^2$ with kernel $U$, so the quotient is $\cong \mathbb{R}^2$ and $\dim U = 2$ (spanned by $x(x - 1)$ and $x^2(x - 1)$). (c) $\ker(T) = \mathrm{Sym}_n$ and $\im(T) = \mathrm{Skew}_n$ (every skew-symmetric $K$ equals $T(K/2)$), so $\mathbb{R}^{n \times n}/\mathrm{Sym}_n \cong \mathrm{Skew}_n$.
    \item[\ref{pr:grassmann}.] (a) $\dim(V \cap W) = 8 - \dim(V + W) \geq 1$, so the possible values are $1, 2, 3, 4$; for example, $\Span\{\mathbf{e}_1, \ldots, \mathbf{e}_4\} \cap \Span\{\mathbf{e}_4, \ldots, \mathbf{e}_7\} = \Span\{\mathbf{e}_4\}$. (b) $U \neq W$, so $U + W = \mathbb{R}^4$ and $\dim(U \cap W) = 3 + 3 - 4 = 2$. Adding and subtracting the two equations gives $x_1 + x_3 = 0$ and $x_2 + x_4 = 0$, so a basis is $(1, 0, -1, 0)^\top$, $(0, 1, 0, -1)^\top$.
    \item[\ref{pr:third-iso}.] $V/U \cong \mathbb{R}^3$ (keep $x_1, x_2, x_3$), $W/U \cong \mathbb{R}$ (keep $x_3$), $V/W \cong \mathbb{R}^2$ (keep $x_1, x_2$); indeed $3 - 1 = 2$.
    \item[\ref{pr:softmax}.] (a) The factor $e^c$ cancels. (b) Both are approximately $(0.090, 0.245, 0.665)$; $e^{1001}$ overflows in floating-point arithmetic, so we compute with the representative $(1001, 1002, 1003) - 1003\cdot\mathbf{1} = (-2, -1, 0)$ instead. (c) Yes: $\frac{\operatorname{softmax}(\mathbf{z})_i}{\operatorname{softmax}(\mathbf{z})_j} = e^{z_i - z_j}$, so equal outputs give $z_i - z_j = w_i - w_j$ for all $i, j$, i.e., $\mathbf{z} - \mathbf{w} \in \Span\{\mathbf{1}\}$.
    \item[\ref{pr:quotient-dual}.] The rows with $a_1 + a_2 + a_3 = 0$, e.g., spanned by $(1, -1, 0)$ and $(1, 0, -1)$.
\end{description}

\section*{Eigenvalues and Eigenvectors}

\begin{description}[leftmargin=0pt, itemsep=0.3em]
    \item[\ref{pr:eig2}.] (a) $p_A(\lambda) = \lambda^2 - 3\lambda - 4 = (\lambda - 4)(\lambda + 1)$; $E_4 = \Span\{(2, 3)^\top\}$ and $E_{-1} = \Span\{(1, -1)^\top\}$. (b) $4 + (-1) = 3$ and $4 \cdot (-1) = -4$. (c) $P = \begin{pmatrix} 2 & 1 \\ 3 & -1 \end{pmatrix}$, $D = \diag(4, -1)$, $P^{-1} = \frac{1}{5}\begin{pmatrix} 1 & 1 \\ 3 & -2 \end{pmatrix}$. (d) $A^k = \frac{1}{5}\begin{pmatrix} 2 \cdot 4^k + 3(-1)^k & 2 \cdot 4^k - 2(-1)^k \\ 3 \cdot 4^k - 3(-1)^k & 3 \cdot 4^k + 2(-1)^k \end{pmatrix}$.
    \item[\ref{pr:eig-triangular}.] The eigenvalues are the diagonal entries $2$, $3$ and $1$. The eigenspaces are spanned by $(1, -1, 5)^\top$ (for $2$), $(0, -2, 1)^\top$ (for $3$) and $(0, 0, 1)^\top$ (for $1$). Diagonalizable, since the three eigenvalues are distinct.
    \item[\ref{pr:eig-defective}.] The eigenvalue $2$ has algebraic multiplicity $2$, but geometric multiplicity $3 - \rank(N - 2I_3) = 1$; the eigenvalue $3$ has both multiplicities equal to $1$. Not diagonalizable.
    \item[\ref{pr:eig-geometric}.] (a) $1$ on the $x$-axis, $-1$ on the $y$-axis. (b) $1$ on $\Span\{(1, 1)\}$, $0$ on $\Span\{(1, -1)\}$. (c) Only $1$, with eigenspace the $x$-axis. (d) No real eigenvalues; the complex eigenvalues are $\pm i$. (e) $2$ on the $x$-axis, $3$ on the $y$-axis. (f) Only $-1$, with eigenspace $\mathbb{R}^2$. All are diagonalizable except (c) and (d) (over $\mathbb{R}$).
    \item[\ref{pr:eig-gre}.] (a) Yes (three distinct eigenvalues); yes ($0$ is not an eigenvalue). (b) $5$, $-8$, $1 + 4 + 16 = 21$, $(-8)^3 = -512$. (c) $A^{-1}$: $-1, \frac{1}{2}, \frac{1}{4}$; $A^2$: $1, 4, 16$; $A + 2I_3$: $1, 4, 6$.
    \item[\ref{pr:eig-idempotent}.] (a) If $P\mathbf{v} = \lambda \mathbf{v}$, then $\lambda \mathbf{v} = P^2 \mathbf{v} = \lambda^2 \mathbf{v}$, so $\lambda^2 = \lambda$. (b) If $\mathbf{v} = Pw$, then $P\mathbf{v} = P^2w = Pw = \mathbf{v}$. Bases of $\im(T_P)$ and $\ker(T_P)$ together span $\mathbb{R}^n$ and are linearly independent, so they form a basis of eigenvectors. (c) Only ``$P$ is diagonalizable''; $P = \begin{pmatrix} 1 & 1 \\ 0 & 0 \end{pmatrix}$ is a counterexample to the other two.
    \item[\ref{pr:eig-props}.] (a) For the shear, $E_1 = \Span\{\mathbf{e}_1\}$, but for its transpose, $E_1 = \Span\{\mathbf{e}_2\}$. (c) $\mathbf{v}^\top \mathbf{v} = (Q\mathbf{v})^\top(Q\mathbf{v}) = \lambda^2 \mathbf{v}^\top \mathbf{v}$. (d) $A^k \mathbf{v} = \lambda^k \mathbf{v} = \mathbf{0}$. (e) $P^{-1} I_2 P = I_2$ for every $P$.
    \item[\ref{pr:eig-recurrence}.] (a) $M = \begin{pmatrix} 1 & 2 \\ 1 & 0 \end{pmatrix}$. (b) Eigenvalues $2$ and $-1$, with eigenvectors $(2, 1)^\top$ and $(-1, 1)^\top$; $a_n = \frac{2^n - (-1)^n}{3}$. (c) $a_{10} = \frac{1024 - 1}{3} = 341$.
    \item[\ref{pr:eig-markov}.] (a) $\lambda = 1$ with eigenvector $(4, 3)^\top$, and $\lambda = 0.3$ with eigenvector $(1, -1)^\top$; the stationary distribution is $\big(\frac{4}{7}, \frac{3}{7}\big)^\top$. (b) $x_k = \big(\frac{4}{7}, \frac{3}{7}\big)^\top + \frac{3}{7}(0.3)^k (1, -1)^\top$; $x_2 = (0.61, 0.39)^\top$; $x_k \to \big(\frac{4}{7}, \frac{3}{7}\big)^\top$.
    \item[\ref{pr:eig-power}.] (a) $(2, 1)$, $(5, 4)$, $(14, 13)$, $(41, 40)$; the ratio tends to $1$, i.e., to the direction of the eigenvector $(1, 1)$. (b) $x_0 = \frac{1}{2}(1, 1)^\top + \frac{1}{2}(1, -1)^\top$, so $A^k x_0 = \frac{1}{2}(3^k + 1, \ 3^k - 1)^\top$.
    \item[\ref{pr:eig-complex}.] (a) $p(\lambda) = (a - \lambda)^2 + b^2$. (b) $3 \pm 4i$; the matrix is $5 R_\theta$ with $\cos\theta = \frac{3}{5}$ and $\sin\theta = \frac{4}{5}$: $r = |3 \pm 4i| = 5$, and $\theta$ is the angle of $3 + 4i$.
    \item[\ref{pr:exp-eigen}.] (a) $(e^{\lambda x})' = \lambda e^{\lambda x}$. If $f' = \lambda f$, then $(f(x)e^{-\lambda x})' = (f'(x) - \lambda f(x))e^{-\lambda x} = 0$, so $f(x)e^{-\lambda x} = c$ is constant, and $f = ce^{\lambda x}$. (b) They are eigenvectors of $D$ for the three different eigenvalues $1$, $2$ and $3$, and the proof of the theorem only uses the linearity of $D$. (c) The functions $e^x, e^{2x}, e^{3x}, \ldots$ are infinitely many linearly independent vectors in $C^\infty(\mathbb{R})$.
\end{description}

\section*{Tensors}

\begin{description}[leftmargin=0pt, itemsep=0.3em]
    \item[\ref{pr:tensor-rules}.] (a) Both are $2(v \otimes w)$. Also, $v \otimes w + v \otimes (-w) = v \otimes (w - w) = v \otimes 0 = 0$, since $v \otimes 0 = v \otimes (0 \cdot 0) = 0\,(v \otimes 0) = 0$. (b) $3\, e_1 \otimes e_1 - e_1 \otimes e_2 + 6\, e_2 \otimes e_1 - 2\, e_2 \otimes e_2 = \left(\begin{smallmatrix} 3 & -1 \\ 6 & -2 \end{smallmatrix}\right)$. (c) The first one is pure: it is $(e_1 + e_2) \otimes (e_1 + e_2)$. The second one is not: its matrix $\left(\begin{smallmatrix} 0 & 1 \\ -1 & 0 \end{smallmatrix}\right)$ has rank $2$, while pure tensors have rank at most $1$. (d) $2 \cdot 3 \cdot 4 = 24$, and $6$, since a dual space has the same dimension.
    \item[\ref{pr:tensor-shapes}.] (a) Order $4$, with the shape $300 \times 3 \times 64 \times 64$ (frames, channels, rows, columns); the space is $\mathbb{R}^{300} \otimes \mathbb{R}^3 \otimes \mathbb{R}^{64} \otimes \mathbb{R}^{64}$, of dimension $300 \cdot 3 \cdot 64 \cdot 64 = 3\,686\,400$. (b) The entry $T_{ijk}$ is $v_{12(i-1) + 4(j-1) + k}$. So $T_{213} = v_{15}$, and $v_{17} = T_{221}$, since $17 = 12 + 4 + 1$.
    \item[\ref{pr:bilinear-practice}.] (a) Bilinear, with $A = \left(\begin{smallmatrix} 2 & -1 \\ -1 & 4 \end{smallmatrix}\right)$. (b) Not bilinear: $\beta(\mathbf{0}, \mathbf{y}) = 1 \neq 0$. (c) Bilinear, with $A = \left(\begin{smallmatrix} 1 & 2 \\ -1 & -2 \end{smallmatrix}\right) = (1, -1)^\top \otimes (1, 2)^\top$. (d) Not bilinear: it is not linear in $\mathbf{x}$; for example, $\beta(2\mathbf{x}, \mathbf{y}) = 4\beta(\mathbf{x}, \mathbf{y})$. Only (a) satisfies $\beta(\mathbf{x}, \mathbf{y}) = \beta(\mathbf{y}, \mathbf{x})$, since a bilinear form is symmetric if and only if its matrix is symmetric.
    \item[\ref{pr:outer-practice}.] (a) $\mathbf{u} \otimes \mathbf{v} = \left(\begin{smallmatrix} 2 & 0 & 1 \\ -2 & 0 & -1 \end{smallmatrix}\right)$, of rank $1$. (b) $T_{ijk} = u_iv_jw_k$, with the slices $w_1(\mathbf{u} \otimes \mathbf{v}) = \left(\begin{smallmatrix} 2 & 0 & 1 \\ -2 & 0 & -1 \end{smallmatrix}\right)$ and $w_2(\mathbf{u} \otimes \mathbf{v}) = \left(\begin{smallmatrix} 6 & 0 & 3 \\ -6 & 0 & -3 \end{smallmatrix}\right)$. (c) No: its determinant is $-2 \neq 0$, so it has rank $2$, while a nonzero outer product has rank $1$. Column by column, $\left(\begin{smallmatrix} 1 & 2 \\ 3 & 4 \end{smallmatrix}\right) = (1, 3)^\top \otimes (1, 0)^\top + (2, 4)^\top \otimes (0, 1)^\top$.
    \item[\ref{pr:einsum-practice}.]
    \begin{enumerate}[label=(\alph*), nosep]
        \item $A_{ij}x_i$, a vector of length $n$: \texttt{np.einsum("ij,i->j", A, x)}.
        \item $A_{ij}B_{ij}$, a number: \texttt{np.einsum("ij,ij->", A, B)}.
        \item $A_{ij}B_{ji}$ with the free index $i$, a vector of length $n$: \texttt{np.einsum("ij,ji->i", A, B)}. It needs $n^2$ multiplications, instead of the $n^3$ for all of $AB$.
        \item $X_{ni}W_{mi}$, a matrix of shape $N \times m$: \texttt{np.einsum("ni,mi->nm", X, W)}.
    \end{enumerate}
    \item[\ref{pr:linear-map-tensor}.] (a) The matrix of $T$ is $A = \left(\begin{smallmatrix} 1 & 2 \\ 0 & 3 \end{smallmatrix}\right)$, so the tensor is $\mathbf{e}_1 \otimes e^1 + 2\, \mathbf{e}_1 \otimes e^2 + 3\, \mathbf{e}_2 \otimes e^2 = \mathbf{e}_1 \otimes (e^1 + 2e^2) + \mathbf{e}_2 \otimes (3e^2)$, one pure tensor for each row of $A$. (b) The contraction is $1 \cdot e^1(\mathbf{e}_1) + 2\, e^2(\mathbf{e}_1) + 3\, e^2(\mathbf{e}_2) = 1 + 0 + 3 = 4 = \tr(A)$. (c) With $P = \left(\begin{smallmatrix} 1 & 1 \\ 1 & -1 \end{smallmatrix}\right)$, $P^{-1}AP = \left(\begin{smallmatrix} 3 & -2 \\ 0 & 1 \end{smallmatrix}\right)$, again with the trace $4$.
    \item[\ref{pr:change-of-basis-practice}.] (a) $P^{-1}AP = \diag(3, -1)$: the columns of $P$ are eigenvectors of $A$. (b) $P^\top AP = \diag(6, -2)$. (c) $A$ and $P^{-1}AP$ both have the trace $2$ and the determinant $-3$, while $P^\top AP$ has the trace $4$ and the determinant $-12 = \det(P)^2\det(A)$. The trace and the determinant are properties of the linear transformation, but not of the bilinear form; for the bilinear form, only the sign of the determinant does not depend on the basis.
    \item[\ref{pr:unfolding-rank}.] (a) The unfolding of $\mathbf{u} \otimes \mathbf{v} \otimes \mathbf{w}$ is the outer product of $\mathbf{u}$ with the vector $(v_jw_k)_{j,k}$ of length $n_2n_3$, and unfolding is linear. So the unfolding of a sum of $r$ rank-one tensors is a sum of $r$ outer products, which has rank at most $r$ by the proposition on matrix rank and outer products. (b) The unfolding of $\mathbf{e}_1 \otimes \mathbf{e}_1 \otimes \mathbf{e}_1 + \mathbf{e}_2 \otimes \mathbf{e}_2 \otimes \mathbf{e}_2$ has the rows $(1, 0, 0, 0)$ and $(0, 0, 0, 1)$, so it has rank $2$. Hence the tensor has rank at least $2$, and it is a sum of two rank-one tensors. (c) $T = (1, 1)^\top \otimes (1, 1)^\top \otimes (1, 1)^\top$.
\end{description}

\end{document}
